\addcontentsline{toc}{chapter}{Abstract}\vspace{-1cm}
%Border
\begin{tikzpicture}[remember picture, overlay]
  \draw[line width = 4pt] ($(current page.north west) + (0.75in,-0.75in)$) rectangle ($(current page.south east) + (-0.75in,0.75in)$);
\end{tikzpicture}


Dehydration remains a critical yet frequently overlooked health issue in everyday environments, especially when exertion, ambient temperature, humidity, and individual physiology jointly influence hydration status. AquaTrack addresses this problem through a prototype-based multimodal health-monitoring system that combines wearable biomarker sensing with context-aware machine learning for real-time dehydration risk prediction.

The proposed system uses a wearable sweat-patch prototype to acquire physiological indicators such as sweat electrolytes, salivary markers, skin measurements, and environmental context. The acquired data is pre-processed using median imputation, winsorization, logarithmic transformation, robust scaling, and domain-informed feature engineering. A dual-pipeline machine learning architecture then trains independent classifiers for the biomarker dataset and the context dataset before combining their outputs through a weighted fusion score.

Experimental observations show that sodium alginate hydrogel patches cross-linked with calcium chloride provide a flexible and functional sensing matrix for sweat collection and urea-responsive colorimetric detection. In the prediction pipeline, the biomarker model and context model achieve strong classification performance, with the fused architecture improving interpretability by identifying cases where physiological risk and contextual risk diverge.

AquaTrack therefore offers a practical framework for mobile health applications in sports science, occupational health, and personal wellness monitoring. By combining non-invasive sensing, hydrogel-based biomarker capture, and lightweight deployment through a FastAPI backend and React Native dashboard, the system demonstrates a reproducible path toward real-time hydration monitoring.

\pagebreak
